\documentclass{report}
\usepackage{amsmath,amssymb}
\setlength{\parindent}{0mm}
\setlength{\parskip}{1em}
\begin{document}
\begin{center}
\rule{6in}{1pt} \
{\large Luis Chacon \\
{\bf An energy- and charge-conserving, fully implicit particle-in-cell algorithm}}

ORNL \\ Oak Ridge \\ TN 37830
\\
{\tt chaconl@ornl.gov}\\
Guangye Chen\\
Daniel C. Barnes\end{center}

Classical particle-in-cell (PIC) algorithms {[}1{]} employ an explicit
time integration approach (leap-frog) to advance the Vlasov-Maxwell
(or Vlasov-Poisson) system using particles coupled to a grid. As in
the fluid context, explicit PIC features a numerical stability constraint
which enforces a minimum timestep resolution (i.e, the timestep has
to be small enough to resolve the fastest time-scale of the system,
e.g., the plasma wave period). Unlike fluid approaches, however, explicit
PIC also features a spatial numerical stability constraint that enforces
a minimum spatial resolution (i.e., the grid size must be smaller
than the Debye length) to avoid the so-called finite-grid instability.
These stringent stability constraints make explicit PIC approaches
impractical for many realistic multi-scale kinetic problems.

Fully implicit temporal schemes can in principle employ arbitrarily
large timesteps without numerical instability. Furthermore, in the
context of PIC methods, they can potentially avoid the finite-grid-instability,
thus allowing spatial resolution to be dictated by accuracy, not stability.
It follows that, by taking time steps and grid sizes much larger than
the plasma period and the Debye length, respectively, implicit PIC
algorithms can be orders of magnitude more efficient than explicit
approaches. This realization motivated the early exploration of various
implicit PIC approaches (namely, moment implicit {[}2,3{]} and direct
implicit {[}4,5{]} methods). However, accurately solving the resulting
nonlinearly coupled system of particle-field equations proved a daunting
task, motivating algorithmic simplifications such as lagging and linearization
which failed to converge the nonlinear residual. As a result, large
numerical errors developed in long-time simulations {[}6{]}, with
significant self-heating/cooling often observed, especially when large
time steps were employed.

The goal of this work is to demonstrate the algorithmic feasibility
and accuracy benefits of a fully implicit PIC algorithm. Our solver
is based on Jacobian-Free Newton-Krylov (JFNK) technology. As a proof-of-principle,
we focus on a one-dimensional electrostatic model {[}7{]} without
preconditioning. One potential show-stopper is the very large nonlinear
residuals that would result if one considered particle quantities
(positions and velocities) as dependent variables. In this case, the
residual size, and hence all vectors in the nonlinear solver, would
scale with the number of particles in the simulation, resulting in
dramatic memory requirements. Our implementation avoids this problem
by eliminating nonlinearly particle quantities in favor of macroscopic
fields and/or moments. Such nonlinear elimination (which we term particle
enslavement) is possible because particle orbits only require macroscopic
fields to evolve, and results in a nonlinear field solver with memory
requirements comparable to those of a fluid computation.

Our implicit implementation is based on the Vlasov-Ampere system.
Our discrete implementation is exactly energy- and charge-conserving,
thus avoiding particle self-heating and self-cooling. While momentum
is not exactly conserved, errors are kept small by an adaptive particle
sub-stepping orbit integrator. Particle moments are orbit averaged
to ensure energy conservation to numerical round-off in the presence
of multiple substeps. As a result, very large time steps, constrained
only by the dynamical time scale of interest, are possible without
appreciable loss of accuracy.

Standard numerical examples are presented that demonstrate the properties
of the scheme. In particular, long-time ion acoustic wave and ion
acoustic shock wave simulations show that numerical accuracy does
not degrade with very large implicit time steps, and that a significant
increase in computational performance is possible versus explicit
simulations, even without a preconditioner. Recent developments to
extend the energy and charge-conserving implicit PIC algorithm to
curvilinear geometry will also be discussed.

\medskip{}


\noindent {\small {[}1{]} C. K. Birdsall, A. B. Langdon, }\textit{\small Plasma
Physics via Computer Simulation,}{\small{} McGraw-Hill, New York, 2004.}{\small \par}

\noindent {\small {[}2{]} R. J. Mason, }\emph{\small J. Comput.
Phys.}{\small ,}\textbf{\small{}
41}{\small{} (2) (1981) 233\textendash{}244.}{\small \par}

\noindent {\small {[}3{]} J. U. Brackbill, D. W. Forslund, }\emph{\small Multiple
time scales,}{\small{} Academic Press, 1985.}{\small \par}

\noindent {\small {[}4{]} A. Friedman, A. B. Langdon, B. I. Cohen,
}\textit{\small Plasma Phys. Controlled Fusion,}\textbf{\small{} 6}{\small{}
(6) (1981) 225 \textendash{} 36.}{\small \par}

\noindent {\small {[}5{]} A. B. Langdon, D. C. Barnes, }\emph{\small Multiple
time scales,}{\small{} Academic Press, 1985.}{\small \par}

\noindent {\small {[}6{]} B. I. Cohen, A. B. Langdon, D. W. Hewett,
R. J. Procassini,}\emph{\small{} J. Comput. Phys.}{\small , }\textbf{\small 81}{\small{}
(1) (1989) 151\textendash{}168. }{\small \par}

\noindent {\small {[}7{]} G. Chen, L. Chac�n, D. C. Barnes, }\emph{\small J.
Comput. Phys.}{\small{} }\textbf{\small 230}{\small , (2011). }{\small \par}

\noindent {\small {[}8{]} G. Chen, L. Chac�n, D. C. Barnes, }\emph{\small J.
Comput. Phys.}{\small{} submitted (2011). }


\end{document}
